\documentclass[11pt]{article}
\usepackage[margin=1.1in]{geometry}
\usepackage{amsmath,amssymb,amsthm,mathtools,booktabs,array,tabularx,longtable}
\usepackage[dvipsnames]{xcolor}
\usepackage[colorlinks=true,linkcolor=MidnightBlue,citecolor=MidnightBlue,urlcolor=MidnightBlue]{hyperref}
\usepackage{fancyhdr}
\pagestyle{fancy}\fancyhf{}\renewcommand{\headrulewidth}{0pt}
\fancyfoot[R]{\thepage}
\fancyfoot[L]{\scriptsize\color{gray}Generated by the Transport Lab from approved parts of its store; not a deposited paper. Papers of record: \href{https://zenodo.org/communities/transport-papers}{Zenodo, Transport Papers}}
\fancypagestyle{plain}{\fancyhf{}\fancyfoot[R]{\thepage}\fancyfoot[L]{\scriptsize\color{gray}Generated by the Transport Lab from approved parts of its store; not a deposited paper. Papers of record: \href{https://zenodo.org/communities/transport-papers}{Zenodo, Transport Papers}}\renewcommand{\headrulewidth}{0pt}}
\newtheorem{theorem}{Theorem}[section]
\newtheorem{proposition}[theorem]{Proposition}
\newtheorem{lemma}[theorem]{Lemma}
\newtheorem{corollary}[theorem]{Corollary}
\newtheorem{conjecture}[theorem]{Conjecture}
\newtheorem{refutation}[theorem]{Refutation}
\newtheorem{prediction}[theorem]{Prediction}
\newtheorem{principle}[theorem]{Principle}
\newtheorem{observation}[theorem]{Observation}
\newtheorem{axiom}[theorem]{Axiom}
\theoremstyle{definition}
\newtheorem{definition}[theorem]{Definition}
\newtheorem{fact}[theorem]{Fact}
\newtheorem{assembled}[theorem]{Assembly}
\newtheorem{postulate}[theorem]{Postulate}
\newtheorem{premise}[theorem]{Premise}
\newtheorem{hypothesis}[theorem]{Hypothesis}
\newtheorem{remark}[theorem]{Remark}
\newtheorem{example}[theorem]{Example}
\theoremstyle{plain}
\setlength{\marginparwidth}{0.8in}
\newcommand{\changedmark}[1]{\marginpar{\raggedright\scriptsize\color{BrickRed}#1}}
\providecommand{\Aa}{}\renewcommand{\Aa}{\mathcal{A}}
\providecommand{\Ad}{}\renewcommand{\Ad}{\widetilde{\mathrm{Ad}}}
\providecommand{\Aev}{}\renewcommand{\Aev}{\mathcal{A}^{+}}
\providecommand{\Aod}{}\renewcommand{\Aod}{\mathcal{A}^{-}}
\providecommand{\B}{}\renewcommand{\B}{\mathrm{B}}
\providecommand{\C}{}\renewcommand{\C}{\mathbb{C}}
\providecommand{\CAR}{}\renewcommand{\CAR}{\mathrm{CAR}}
\providecommand{\CS}{}\renewcommand{\CS}{{\mathbb{S}^{+}}}
\providecommand{\CSm}{}\renewcommand{\CSm}{{\mathbb{S}^{-}}}
\providecommand{\Cl}{}\renewcommand{\Cl}{\mathrm{Cl}}
\providecommand{\Con}{}\renewcommand{\Con}{\mathcal{C}}
\providecommand{\Dl}{}\renewcommand{\Dl}{\Delta}
\providecommand{\EN}{}\renewcommand{\EN}{E_{N}}
\providecommand{\End}{}\renewcommand{\End}{\mathrm{End}}
\providecommand{\F}{}\renewcommand{\F}{\mathbb{F}}
\providecommand{\G}{}\renewcommand{\G}{\Gamma}
\providecommand{\HH}{}\renewcommand{\HH}{\mathbb{H}}
\providecommand{\IS}{}\renewcommand{\IS}{I\sigma_3}
\providecommand{\J}{}\renewcommand{\J}{\mathbf{J}}
\providecommand{\LH}{}\renewcommand{\LH}{L_{H}}
\providecommand{\Mat}{}\renewcommand{\Mat}{\mathrm{M}}
\providecommand{\Mo}{}\renewcommand{\Mo}{M^{\circ}}
\providecommand{\Mp}{}\renewcommand{\Mp}{M^{+}}
\providecommand{\Msun}{}\renewcommand{\Msun}{M_{\odot}}
\providecommand{\Mt}{}\renewcommand{\Mt}{\widetilde{M}}
\providecommand{\Ncone}{}\renewcommand{\Ncone}{\mathcal{N}}
\providecommand{\Neff}{}\renewcommand{\Neff}{N_{\mathrm{eff}}}
\providecommand{\PCS}{}\renewcommand{\PCS}{\mathbb{P}{\mathbb{S}^{+}}}
\providecommand{\PT}{}\renewcommand{\PT}{\mathbb{PT}}
\providecommand{\Pin}{}\renewcommand{\Pin}{\mathrm{Pin}}
\providecommand{\Q}{}\renewcommand{\Q}{\mathbb{Q}}
\providecommand{\R}{}\renewcommand{\R}{\mathbb{R}}
\providecommand{\Ric}{}\renewcommand{\Ric}{\mathrm{Ric}}
\providecommand{\SBH}{}\renewcommand{\SBH}{S_{\mathrm{BH}}}
\providecommand{\Sc}{}\renewcommand{\Sc}{\mathbb{S}}
\providecommand{\Sch}{}\renewcommand{\Sch}{\mathrm{P}}
\providecommand{\Sig}{}\renewcommand{\Sig}{\Sigma}
\providecommand{\Som}{}\renewcommand{\Som}{S_\omega}
\providecommand{\Sp}{}\renewcommand{\Sp}{\mathbb{S}}
\providecommand{\Sph}{}\renewcommand{\Sph}{\mathrm{S}}
\providecommand{\Spi}{}\renewcommand{\Spi}{S_\pi}
\providecommand{\Spin}{}\renewcommand{\Spin}{\mathrm{Spin}}
\providecommand{\Sth}{}\renewcommand{\Sth}{S_{\mathrm{th}}}
\providecommand{\T}{}\renewcommand{\T}{\mathbb{T}}
\providecommand{\TH}{}\renewcommand{\TH}{T_{\mathrm{H}}}
\providecommand{\Tr}{}\renewcommand{\Tr}{\operatorname{Tr}}
\providecommand{\UHF}{}\renewcommand{\UHF}{\mathrm{UHF}}
\providecommand{\Z}{}\renewcommand{\Z}{\mathbb{Z}}
\providecommand{\Zc}{}\renewcommand{\Zc}{\mathrm{Z}}
\providecommand{\cmg}{}\renewcommand{\cmg}{\mathrm{cm^{2}\,g^{-1}}}
\providecommand{\dS}{}\renewcommand{\dS}{\mathrm{dS}}
\providecommand{\dsc}{}\renewcommand{\dsc}{\mathrm{dsc}}
\providecommand{\dtPG}{}\renewcommand{\dtPG}{dt}
\providecommand{\dtS}{}\renewcommand{\dtS}{dt_{\mathrm{S}}}
\providecommand{\ellv}{}\renewcommand{\ellv}{\ell}
\providecommand{\er}{}\renewcommand{\er}{e_r}
\providecommand{\et}{}\renewcommand{\et}{e_t}
\providecommand{\ev}{}\renewcommand{\ev}{\mathrm{ev}}
\providecommand{\gl}{}\renewcommand{\gl}{\mathfrak{gl}}
\providecommand{\hatV}{}\renewcommand{\hatV}{\widehat{V}}
\providecommand{\hb}{}\renewcommand{\hb}{\bar{h}}
\providecommand{\hf}{}\renewcommand{\hf}{\hat{f}}
\providecommand{\id}{}\renewcommand{\id}{\mathrm{id}}
\providecommand{\kpc}{}\renewcommand{\kpc}{\mathrm{kpc}}
\providecommand{\lP}{}\renewcommand{\lP}{\ell_{P}}
\providecommand{\mD}{}\renewcommand{\mD}{\mathbf{m}_{D}}
\providecommand{\mb}{}\renewcommand{\mb}{\bar m}
\providecommand{\mnu}{}\renewcommand{\mnu}{\Sigma m_{\nu}}
\providecommand{\nInf}{}\renewcommand{\nInf}{n_\infty}
\providecommand{\nO}{}\renewcommand{\nO}{n_0}
\providecommand{\ninf}{}\renewcommand{\ninf}{n_{\infty}}
\providecommand{\nz}{}\renewcommand{\nz}{n_{0}}
\providecommand{\proj}{}\renewcommand{\proj}{\operatorname{proj}}
\providecommand{\pv}{}\renewcommand{\pv}{\pi}
\providecommand{\rad}{}\renewcommand{\rad}{\mathrm{rad}}
\providecommand{\sgn}{}\renewcommand{\sgn}{\operatorname{sgn}}
\providecommand{\shA}{}\renewcommand{\shA}{\mathcal{A}}
\providecommand{\shC}{}\renewcommand{\shC}{\mathcal{C}}
\providecommand{\shE}{}\renewcommand{\shE}{\mathcal{E}}
\providecommand{\shJ}{}\renewcommand{\shJ}{\mathcal{J}}
\providecommand{\shK}{}\renewcommand{\shK}{\mathcal{K}}
\providecommand{\shL}{}\renewcommand{\shL}{L}
\providecommand{\shW}{}\renewcommand{\shW}{\mathcal{W}}
\providecommand{\sll}{}\renewcommand{\sll}{\mathfrak{sl}}
\providecommand{\so}{}\renewcommand{\so}{\mathfrak{so}}
\providecommand{\su}{}\renewcommand{\su}{\mathfrak{su}}
\providecommand{\supp}{}\renewcommand{\supp}{\operatorname{supp}}
\providecommand{\tr}{}\renewcommand{\tr}{\operatorname{tr}}
\providecommand{\uu}{}\renewcommand{\uu}{\mathfrak{u}}
\providecommand{\br}{}\renewcommand{\br}[1]{\langle #1\rangle}
\providecommand{\bra}{}\renewcommand{\bra}[1]{\langle #1|}
\providecommand{\file}{}\renewcommand{\file}[1]{\texttt{#1}}
\providecommand{\hp}{}\renewcommand{\hp}[2]{( #1 \,|\, #2 )}
\providecommand{\ip}{}\renewcommand{\ip}[2]{\langle #1, #2\rangle}
\providecommand{\ket}{}\renewcommand{\ket}[1]{|#1\rangle}
\providecommand{\rel}{}\renewcommand{\rel}[2]{D\bigl(#1 \,\|\, #2\bigr)}
\providecommand{\rev}{}\renewcommand{\rev}[1]{\widetilde{#1}}
\providecommand{\slashed}{}\renewcommand{\slashed}[1]{{#1}\mkern-10mu/}
\title{\bfseries What Fixes the Angle: The spectrum-condition route to the Hawking temperature, run in the ascent's own algebra: Companion note to T5: within the supplied dynamics positivity of energy selects the Boulware state; a hop across the horizon is what the route needs, and the hop with the sea filled is the pairing}
\author{Leon Butler\\ \normalsize\href{mailto:lbutler.math@gmail.com}{\texttt{lbutler.math@gmail.com}}\\ \normalsize ORCID: \href{https://orcid.org/0009-0006-4449-5617}{0009-0006-4449-5617}}
\date{Generated by the Transport Lab, 27 October 2026}
\begin{document}
\maketitle
\begin{abstract}
\normalsize
The singularity paper fixes the Hawking temperature by a premise, regularity at the fixed locus of the Killing boost, proves that nothing inside the mode algebra replaces it, and records one route to a better premise: positivity of the energy of the boundary modes for a generator the ascent installs, from which the Bisognano--Wichmann mechanism derives the factor $2\pi$ rather than assuming it. This note runs that route as far as the algebra goes, in exact arithmetic, with kill conditions fixed in advance. The geometric half of the mechanism is in the ascent: the boost rotor at imaginary rapidity $i\pi$ is the reflection through the bifurcation surface, and the conformal completion installs the translations as null rotations fixing the point at infinity. The selective half is not. Within the supplied dynamics every generator is diagonal in the joint number basis, so a unique ground state is a product across every cut, Gibbs at no finite inverse temperature: the spectrum condition selects a Boulware-type state or nothing. A hop across the cut changes this: the generator $7N_{0} - 12(a^{\dagger}_{0}a_{1} + a^{\dagger}_{1}a_{0})$ has a unique entangled ground state whose one-sided marginal is the Gibbs state at the singularity paper's sample point, and the particle-hole map on one mode carries it to that paper's thermofield pair. The route therefore survives only with the kinetic term the singularity paper's locality premise excludes, its two open items becoming one; the value $2\pi$ remains the imported theorem's, being a ratio of couplings in the finite model.
\end{abstract}
\medskip
\noindent\textbf{MSC 2020.} Primary 83C57; Secondary 15A66, 46L60, 81T05, 82B10.\\
\textbf{Keywords.} Hawking temperature; spectrum condition; Bisognano--Wichmann theorem; KMS state; Boulware state; particle-hole symmetry; conformal model; transport principle.
\newpage
\noindent \textbf{Authorship and disclosure. }The founding idea of the series is the author's: that division by zero in a Clifford algebra should be read not as a prohibition but as a transport into a higher-dimensional algebra, and that the transported division need not be accompanied by an inverse operation. The direction of investigation and the critical review of the results are likewise the author's. The mathematical development and the drafting of the text were carried out in extended collaboration with Claude Fable 5 (Anthropic), an AI system.

\noindent\fbox{\parbox{\dimexpr\linewidth-2\fboxsep-2\fboxrule}{\textbf{This is not a deposited paper.} It was generated by the Transport Lab on 27 October 2026 and assembled from approved parts of its store. It follows results across several papers of the series; each is stated with its proof and credited to the paper that first proved it. The official versions of those papers are on Zenodo, in the Transport Papers community, \url{https://zenodo.org/communities/transport-papers}.}}
\newpage

\section{Setting, and what is imported}

\phantomsection\label{c:note.angle.setting-setting-what-imported}
The singularity paper obtains the Hawking temperature in three moves. Modular descent gives the form: a pure, boost-invariant, entangling state of a horizon pair descends to a Gibbs state for the boost at some inverse temperature $\beta$ in rapidity \cite[Thm.~7.3]{ButlerT5}. Rate rigidity shows $\beta$ is a ratio of two clocks, the state's modular clock and the boost's geometric clock, and that no automorphism of the mode algebra fixes it \cite[Thm.~7.16]{ButlerT5}. Premise (R) then fixes it: each single-mode marginal is KMS for the Killing flow at the period $2\pi/\kappa$ that the smoothness of the Euclidean section dictates \cite[Prop.~7.18, Prem.~7.19]{ButlerT5}, whence $T = \kappa/2\pi = 1/8\pi M$. The paper proves that (R) cannot be internal: the locus where the calibration happens, the bifurcation surface, is one where division does not fail, so the transport principle makes no statement there \cite[Prop.~7.25]{ButlerT5}. It lists the standard forms of the same import, Euclidean regularity, regularity on the future horizon, nonsingularity of the quasi-free state, and the short-distance condition near the horizon, and it records one route to a better-placed premise: the Bisognano--Wichmann theorem does not assume the factor $2\pi$ but derives it from a spectrum condition, positivity of the energy, together with covariance and locality; the ascent to $\Cl(2,4)$ installs translations in the parabolic stabiliser of the point at infinity; so the premise could be restated as positivity of the energy of the boundary modes with respect to a generator the ascent itself installed. Transcribing the theorem, the paper says, needs an action of the relevant subgroup on the boundary modes, a wedge and locality structure, and analyticity of the state, each of which could fail \cite[Rem.~7.26, Rem.~7.33(a)]{ButlerT5}.

This note runs that route as far as exact arithmetic reaches, and no further. What it can test is whether, within the dynamics the singularity paper supplies, a spectrum condition selects a state at all, and what the selected state is; what the ascent contributes to the geometric half of the mechanism; and what a generator outside the supplied class does. What it cannot test is the analytic half, the theorem that turns positivity of energy into the number $2\pi$, which is a statement about infinitely many modes and analytic continuation and is imported below with its hypotheses. The kill conditions of the section ``The claim, and the kill conditions fixed before computation'' were fixed before the computation, after reading the passages of the singularity paper just cited.

Vocabulary, for a reader of the singularity paper: the \emph{supplied class} is the set of generators its completion provides for the boundary-mode dynamics, the second quantisations of function-linear maps of the mode line bundle and the products of the parity sections, all diagonal in the joint number basis, so that no supplied generator moves occupation from one point to another \cite[Prem.~7.30, Thm.~C.8]{ButlerT5}; a \emph{hop} is a generator $a^{\dagger}_{x}a_{y} + a^{\dagger}_{y}a_{x}$ between two modes, the second quantisation of a first-order differential operator, which the supplied class does not contain and whose presence in nature the paper declares as its residual question \cite[Rem.~7.33(b), Rem.~C.9]{ButlerT5}; a \emph{cut} divides a set of modes into two sides, the two ends of a horizon pair; a state is \emph{boost-KMS at $\beta$} on a side if its marginal there is $e^{-\beta K}/Z$ for the boost generator $K$ restricted to that side; the \emph{spectrum condition} for a generator $H$ is that the state is the unique ground state of $H$; and the \emph{Boulware} member of the singularity paper's dictionary is the product state with pure marginals, $\beta = \infty$ \cite[Rem.~7.26]{ButlerT5}.

\begin{assembled}[Imported: the singularity paper's state selection]\label{c:angle.imported-singularity-papers-state-selection}
For the horizon pair, a state vector that is invariant up to phase under the boost and cyclic and separating for the pair factor has the form $\cos\theta\,\ket{00} + \sin\theta\,\ket{11}$ with $\cos\theta\sin\theta \neq 0$; its one-sided marginal is $\operatorname{diag}(\cos^{2}\theta, \sin^{2}\theta)$, the Gibbs state for the number operator at $\tan^{2}\theta = e^{-\beta\omega}$; the sample point $e^{-\omega/T} = \tfrac{9}{16}$ has $(\cos\theta, \sin\theta) = (\tfrac45, \tfrac35)$. Every automorphism of the mode factor intertwining the boost's phase flow $\operatorname{Ad}(e^{itK})$ acts on $\beta$ through $\{\pm 1\}$ only. The surface gravity is $\kappa = 1/4M$, the Euclidean section is smooth at the bifurcation sphere exactly for the period $2\pi/\kappa$, and (R) imposes KMS at that period mode by mode. The supplied class is ultralocal, and the s-wave array of Step 4e has weights $\omega_{k} = \kappa k$ with its sea level at $k = 0$. \emph{Provenance:} \cite[Lem.~7.2, Thm.~7.3, Thm.~7.16, Prop.~7.18, Prem.~7.19, Prop.~7.25, Rem.~7.26, Prem.~7.30, Rem.~7.33, Prem.~8.1, Thm.~C.8, Rem.~C.9, Prop.~A.2]{ButlerT5}.
\end{assembled}
\begin{proof}
For the horizon pair, a state vector that is invariant up to phase under the boost and cyclic and separating for the pair factor has the form $\cos\theta\,\ket{00} + \sin\theta\,\ket{11}$ with $\cos\theta\sin\theta \neq 0$ (the cyclic and separating hypothesis excludes the product vectors $\ket{01}$, $\ket{10}$, which are also invariant up to phase); its one-sided marginal is $\operatorname{diag}(\cos^{2}\theta, \sin^{2}\theta)$, the Gibbs state for the number operator at $\tan^{2}\theta = e^{-\beta\omega}$; the sample point $e^{-\omega/T} = \tfrac{9}{16}$ has $(\cos\theta, \sin\theta) = (\tfrac45, \tfrac35)$. Every automorphism of the mode factor intertwining the boost's phase flow $\operatorname{Ad}(e^{itK})$ acts on $\beta$ through $\{\pm 1\}$ only. The surface gravity is $\kappa = 1/4M$, the Euclidean section is smooth at the bifurcation sphere exactly for the period $2\pi/\kappa$, and (R, \ref{c:singularity.r-regularity-fixed-locus}) imposes KMS at that period mode by mode. The supplied class (\ref{c:singularity.l-boundary-mode-locality}) is ultralocal, and the s-wave array of Step 4e has weights $\omega_{k} = \kappa k$ with its sea level at $k = 0$ (by \ref{c:singularity.q-joint-quasi-freeness-s-wave}). \emph{Provenance:} \cite[Lem.~7.2, Thm.~7.3, Thm.~7.16, Prop.~7.18, Prem.~7.19, Prop.~7.25, Rem.~7.26, Prem.~7.30, Rem.~7.33, Prem.~8.1, Thm.~C.8, Rem.~C.9, Prop.~A.2]{ButlerT5} (\ref{c:singularity.rate-rigidity-temperature-not-algebra-intrinsic}, \ref{c:singularity.modular-descent}, \ref{c:singularity.schmidt-forcing-proved-outright}, \ref{c:singularity.euclidean-section-period}).

This restates \ref{c:singularity.class-relative-no-go-no-supplied}, \ref{c:singularity.computed}.

\end{proof}

\begin{assembled}[Imported: the conformal model and the null rotations]\label{c:angle.imported-conformal-model-null-rotations}
The completion $\Cl(p{+}1,q{+}1)$ of $\Cl(p,q)$ carries the null pair $\nO = \tfrac12(e_{-} - e_{+})$, $\nInf = e_{-} + e_{+}$ with $\nO \cdot \nInf = -1$, and represents the point $x$ by the null vector $X(x) = \nO + x + \tfrac12 Q(x)\nInf$, with $X(x)\cdot X(y) = -\tfrac12 Q(x - y)$. For a null vector $v$ and $u$ orthogonal to it, $R_{u} = 1 + vu$ is a rotor fixing $v$; and for a null partner $w$ of $v$ ($v \cdot w \neq 0$) and $U = \{v, w\}^{\perp}$, the nondegenerate complement of the hyperbolic plane they span, $u \mapsto R_{u}$ (for $u \in U$) is an isomorphism of $(U, +)$ onto the group of null rotations of $v$. \emph{Provenance:} \cite[Sec.~4.3, Lem.~4.6, Thm.~4.10]{ButlerT1}.
\end{assembled}
\begin{proof}
The completion $\Cl(p{+}1,q{+}1)$ of $\Cl(p,q)$ carries the null pair $\nO = \tfrac12(e_{-} - e_{+})$, $\nInf = e_{-} + e_{+}$ with $\nO \cdot \nInf = -1$, and represents the point $x$ by the null vector $X(x) = \nO + x + \tfrac12 Q(x)\nInf$, with $X(x)\cdot X(y) = -\tfrac12 Q(x - y)$. For a null vector $v$ and $u$ orthogonal to it, $R_{u} = 1 + vu$ is a rotor fixing $v$, and $u \mapsto R_{u}$ is an isomorphism of $(U, +)$ onto the group of null rotations of $v$. \emph{Provenance:} \cite[Sec.~4.3, Lem.~4.6, Thm.~4.10]{ButlerT1} (\ref{c:transport.group-null-rotations}, \ref{c:transport.totality}).

\end{proof}

\begin{fact}[Imported: the theorems that deliver $2\pi$, and their hypotheses]\label{c:angle.imported-theorems-deliver-their-hypotheses}
For a Wightman field theory in Minkowski space, the modular group of the algebra of a wedge with respect to the vacuum is the one-parameter group of boosts preserving the wedge, with modular parameter $t$ corresponding to rapidity $2\pi t$, and the modular conjugation is the CPT operator composed with the rotation by $\pi$ about the boost axis, which acts on spacetime as the wedge reflection $(x^{0}, x^{1}) \mapsto (-x^{0}, -x^{1})$; the hypotheses are the spectrum condition (positivity of the energy-momentum), Poincar\'e covariance and locality. The abstract form: if a von Neumann algebra $M$ with cyclic separating vector $\Omega$ carries a one-parameter unitary group $U(a)$ with positive generator, $U(a)\Omega = \Omega$ and $U(a)MU(a)^{*} \subseteq M$ for $a \geq 0$, then the modular group of $(M, \Omega)$ acts on the translations as the dilation $\Delta^{it}U(a)\Delta^{-it} = U(e^{-2\pi t}a)$. For the free massless field the modular groups of the wedge and double-cone algebras are geometric, and for a conformal net on the circle the modular group of an interval's algebra is the M\"obius flow fixing the interval's endpoints, again at $2\pi$. On a spacetime with a bifurcate Killing horizon the same mechanism gives the Hawking temperature for states regular across the horizon, and a stationary quasi-free state that is regular is unique and KMS at $\kappa/2\pi$. \emph{Provenance:} \cite{BW1975, BW1976}; \cite{Borchers1992}; \cite{HislopLongo1982, BGL1993}; \cite{Sewell1982, KayWald1991}; the Rindler case \cite{Unruh1976}.
\emph{Source:} J. J. Bisognano and E. H. Wichmann, On the duality condition for a Hermitian scalar field, J. Math. Phys. 16 (1975), 985--1007.; J. J. Bisognano and E. H. Wichmann, On the duality condition for quantum fields, J. Math. Phys. 17 (1976), 303--321.; H.-J. Borchers, The CPT-theorem in two-dimensional theories of local observables, Commun. Math. Phys. 143 (1992), 315--332.; P. D. Hislop and R. Longo, Modular structure of the local algebras associated with the free massless scalar field theory, Commun. Math. Phys. 84 (1982), 71--85.; R. Brunetti, D. Guido and R. Longo, Modular structure and duality in conformal quantum field theory, Commun. Math. Phys. 156 (1993), 201--219.; G. L. Sewell, Quantum fields on manifolds: PCT and gravitationally induced thermal states, Ann. Phys. 141 (1982), 201--224.; B. S. Kay and R. M. Wald, Theorems on the uniqueness and thermal properties of stationary, nonsingular, quasifree states on spacetimes with a bifurcate Killing horizon, Phys. Rep. 207 (1991), 49--136.; W. G. Unruh, Notes on black-hole evaporation, Phys. Rev. D 14 (1976), 870--892..
\end{fact}

\phantomsection\label{c:note.angle.setting-setting-what-imported-2}
The abstract form is the one the route needs: a half-sided translation group with positive generator on the mode algebra, leaving the state invariant, forces the modular group to act on the translations as the dilation by $e^{-2\pi t}$. What the theorem asks of the framework is therefore three things: a translation group the ascent installs, a state with positive energy for it that is unique and cyclic and separating for one side, and the half-sided inclusion, the locality structure. The note tests the first in the geometry and a two-mode analogue of the second, and identifies the rest as the analytic import.

\section*{What we assume}
The results below are taken as given in this paper; each result that rests on one says so.
\begin{postulate}[(L$_E$): boundary-mode locality, as descent covariance on finite windows]\label{c:singularity.l-boundary-mode-locality}
\emph{Premise (L$_{E}$), boundary-mode locality.} On every finite window of the boundary modes, $v_{1}, \dots, v_{n}$ satisfying the CAR, the internal dynamics $\alpha_{t}$ of the boundary modes commutes with the descent at all times:
\[
\alpha_{t}\circ\mathcal{E} \;=\; \mathcal{E}\circ\alpha_{t} \quad \text{for all } t, \qquad \mathcal{E} = E_{1}\circ\cdots\circ E_{n}, \quad E_{k}(x) = N_{k}xN_{k} + (1 - N_{k})x(1 - N_{k}),
\]
$\mathcal{E}$ the descent of the window (Proposition \emph{The descent, and what it erases} [T5, Prop 5.5]).

\emph{The version 1 wording, equivalent.} By Theorem \emph{On finite windows, boundary-mode locality is equivalent to the dynamics commuting with the descent for all times}, on every finite window this holds if and only if the dynamics is generated by elements of the abelian algebra $\mathcal{D}$ that the $\sigma_{k} = 2N_{k} - 1$ generate, diagonal in the joint number basis: the internal dynamics is generated by the geometric data the framework supplies on $E$, the boost bivector field $x \mapsto \sigma_{r}(x)$ and its products. Cross-point products such as $\sigma_{j}\sigma_{k}$ correlate phases without moving occupation; hopping $v_{j}^{*}v_{k} + v_{k}^{*}v_{j}$, pairing $v_{j}^{*}v_{k}^{*} + v_{k}v_{j}$ and every other generator with a component outside $\mathcal{D}$ are excluded by the one clause. Within the class of generators the framework supplies, the exclusion of hopping is a theorem (Theorem \ref{c:singularity.class-relative-no-go-no-supplied}); what the premise adds is the exclusion of unsupplied generators, pairing terms and the possible differential dynamics, which Proposition \emph{(L$_E$) as the stabiliser and the tangent} [T5, Prop 13.4] places in the tangent of the vacuum orbit and in its stabiliser.

\emph{What it forbids.} By Theorem \emph{Number-diagonal dynamics carries operator support only to directly coupled modes, with no exponential growth, so no dynamics (L\_E) admits realises the kappa\^{}{-1} ln N spreading law in real time}, every dynamics the premise admits carries an evolved mode operator only to the modes directly coupled to it, and every local commutator and out-of-time-order correlator is almost periodic and bounded: the premise forbids exponential real-time commutator growth. The real-time spreading law $t_{*} = \kappa^{-1}\ln N$ of Proposition \emph{The spreading law, and the chaos bound met, not saturated} [T8, Asm 6.1] is therefore not realised under it, and the scrambling time is a modular-time statement: the growth $e^{\kappa t}$ is a rate of the rotor flow, the KMS-strip continuation of the phase flow that is the windows' time evolution (Theorem \emph{The phase flow is the time evolution; the rotor flow is its KMS continuation}). The premise is kept over fast real-time scrambling; the two cannot both hold.
\end{postulate}

\begin{postulate}[(Q): joint quasi-freeness, with the s-wave weights at the half-integer offset]\label{c:singularity.q-joint-quasi-freeness-s-wave}
In the stationary idealisation of the boundary modes, and for this step alone, the joint state of the mode array is taken to be the gauge-invariant quasi-free state with the marginal data that (P) and (R) (Premise \ref{c:singularity.r-regularity-fixed-locus}) fix unconditionally (Theorem \emph{The rate, and the discharge of hypothesis (iii)} [T5, Thm 7.20], Corollary \emph{The chain's state class, and its stability} [T5, Cor 7.32]): writing $\psi_k$ for the fermion modes of the array of Theorem \emph{The bosonisation bridge, s-wave sector} [T5, Thm 8.2], its two-point function is $\langle \psi^\dagger_k \psi_l \rangle = \delta_{kl}\,n_k$, and every higher correlation function factorises by Wick's theorem, so that all connected correlators vanish and the joint state is KMS at $\beta = 2\pi/\kappa$ for the joint modular flow. The mode weights are evenly spaced with the array's sea level midway between two sites,
\[
\omega_k = \kappa\bigl(k + \tfrac12\bigr), \qquad k \in \Z ,
\]
so that no mode has zero weight, the weights are symmetric under the particle-hole relabelling $k \mapsto -1-k$, and the Killing flow at one modular period $2\pi/\kappa$ is the parity automorphism $\psi \mapsto -\psi$.

\emph{Why this offset.} The weights $\omega_k$ of Definition \emph{The geometric modular condition (G)} [T5, Def 7.12] are fixed nowhere else in the framework: given (P), (R) and the quasi-free clause, an affine law $\omega_k = ck + d$ is what the currents' being normal modes of the boost flow requires, particle-hole symmetry leaves exactly two offset classes, $d = 0$ and $d = c/2$, and the slope is the value $\beta c = 2\pi$ (Theorem \emph{Premise (Q) splits into the maximum-entropy member of the (P),(R) class and an affine weight law with two free data}(ii), (iv)). The class $d = c/2$ is adopted as a consequence of one reading, not as a bare premise: the array's modes are read as the discretised modes of the field theory's spinor fields on the horizon's light rays, $\psi(u) \mapsto e^{\kappa t/2}\psi(e^{\kappa t}u)$, of boost weight $\tfrac12$ (the weight at which the boost preserves the anticommutator), in the state that Premise \emph{(N$_{+}$): on the continuum of a horizon's null generators the state of the modes is a gauge-invariant quasi-free ground state for the affine null translations along the generators} (N$_{+}$) selects, a ground state for the affine null translations of the generators (Theorem \emph{A ground state for the affine null translations of a horizon's generators restricts on either side of the edge to the continuum I-state at rapidity period $2\pi$, so the continuum restatements of (R) for the eternal hole, (R$_{\Lambda}$), (R$_{\mathrm{loc}}$) and the joint KMS part of (Q) follow from one premise with no temperature in it, posable only on the continuum}). A field of boost weight $\tfrac12$ changes sign under a full turn of the regular Euclidean section, and $\kappa(k + \tfrac12)$ is the offset at which the array's Killing flow at one modular period carries that sign (Theorem \emph{With the sea level between sites the s-wave array carries exactly half of ln 2 at the Hawking temperature}). The reading is a decision, not a derivation: the Euclidean-period sign $R_{2\pi} = -1$ (Proposition \emph{Where the $+1$ of Fermi-Dirac comes from} [T5, Cor 5.4]) holds in a KMS state of the array at either offset, so identifying the array's real-time recurrence at one period with it is part of the reading.

\emph{On the continuum.} The joint KMS part of this premise is not a separate assumption on the continuum mode space of a horizon's generators: there it follows from (N$_{+}$), whose null ground state restricts on either side of the edge to the continuum I-state at $\beta = 2\pi/\kappa$, jointly KMS and even (Theorem \emph{A ground state for the affine null translations of a horizon's generators restricts on either side of the edge to the continuum I-state at rapidity period $2\pi$, so the continuum restatements of (R) for the eternal hole, (R$_{\Lambda}$), (R$_{\mathrm{loc}}$) and the joint KMS part of (Q) follow from one premise with no temperature in it, posable only on the continuum}(ii), (iii)). The stationary array statement here is its point-spectrum discretisation, on which (N$_{+}$) itself cannot be posed (part (iv) there). The quasi-free clause is a hypothesis of (N$_{+}$) as well, and (N$_{+}$) does not give the weight clause.

\emph{Consequences of the offset.} At these weights one array carries exactly $\tfrac12\ln 2$ at the Hawking temperature, given the classical theta-function identities, so under the entropy-matching hypothesis two arrays carry the capacity $\ln 2$ of one registered boundary mode with no zero-weight mode (Theorem \emph{With the sea level between sites the s-wave array carries exactly half of ln 2 at the Hawking temperature}(i), (iii)); the current sector of one array is one oscillator tower, one helicity channel, at either offset (Theorem \emph{Under (Q) an array's ln 2 is its zero-mode charge doublet, which no current with m != 0 connects, and its current sector is one oscillator tower}(ii), (iv)), so one registered mode then corresponds to two arrays and two helicity channels. Since every mode has nonzero weight, the gluing equation of Theorem \emph{The rate, and the discharge of hypothesis (iii)} [T5, Thm 7.20] fixes $\beta$ at every mode of the array, and the zero-weight scopes of Theorems \emph{The descended form is state independent} [T5, Thm 7.6], \emph{Universality of the temperature under (G)} [T5, Cor 7.15] and \emph{The rate, and the discharge of hypothesis (iii)} [T5, Thm 7.20] do not arise on it.

\emph{Version 2.} The source states the weight clause as $\omega_k = \kappa k$ with the sea level at $k = 0$, one mode of zero weight at which the gluing equation fixes $n_0 = \tfrac12$ but not $\beta$. Results computed under that offset refer to version 2.
\end{postulate}

\begin{postulate}[(R): regularity at the fixed locus]\label{c:singularity.r-regularity-fixed-locus}
For every mode $k$ and every pair $a, b$ of elements of its ladder $\{v_k, v_k^{*}\}$, the correlation function $F_{a,b}(t) := \omega\bigl(a\,\alpha_t(b)\bigr)$ of the Killing flow $\alpha_t = \operatorname{Ad}(e^{it\omega_k K_k})$, an entire function of $t$, continues in imaginary Killing time to a single-valued function on the regular Euclidean section of the Schwarzschild gravitational field on flat spacetime (Proposition \ref{c:singularity.euclidean-section-period}; Definition \emph{Gauge Theory Gravity: the gravitational field on flat spacetime}): its two orderings are the two sides of the cut, and they meet after one period,
\[
F_{a,b}\bigl(t + i\,p\bigr) \;=\; \omega\bigl(\alpha_t(b)\,a\bigr), \qquad p = \frac{2\pi}{\kappa} .
\]
In words: each single-mode marginal is KMS for the Killing flow at exactly the period the Euclidean section's smoothness at its tip dictates. The regularity is local, posed at the tip of the Euclidean section (a polar chart about the tip of the Wick-rotated flat background); for a state that is KMS for the Killing flow, this is also the period at which it is regular across the horizon locus in the Painlev\'e-Gullstrand chart.

\emph{Note (continuum restatement).} Posed on the continuum of the horizon's null generators (Definition \emph{Continuum mode space} [T7, Def 4.1]), the content of (R) for the eternal (Kruskal) hole follows from Premise \emph{(N$_{+}$): on the continuum of a horizon's null generators the state of the modes is a gauge-invariant quasi-free ground state for the affine null translations along the generators}: by Theorem \emph{A ground state for the affine null translations of a horizon's generators restricts on either side of the edge to the continuum I-state at rapidity period $2\pi$, so the continuum restatements of (R) for the eternal hole, (R$_{\Lambda}$), (R$_{\mathrm{loc}}$) and the joint KMS part of (Q) follow from one premise with no temperature in it, posable only on the continuum}(iii) the generators' modes are then jointly KMS for the Killing flow at $\beta = 8\pi M$. The wording above, on the ladders $\{v_{k}, v_{k}^{*}\}$ with a phase-acting flow, is a discretisation of that statement, recovered only for modes of vanishing Mellin width; (N$_{+}$) itself cannot be posed on point spectrum. A hole formed by collapse is not covered by (N$_{+}$), and (R) remains a premise. (The whole generator of the eternal horizon, which that case of (N$_{+}$) uses, lies beyond one flat chart; see that premise.)
\end{postulate}

\section{The claim, and the kill conditions fixed before computation}

\phantomsection\label{c:note.angle.setting-claim-kill-conditions-fixed-before}
The kill conditions and predictions below are as registered \cite{ButlerPreregT05angle}; every departure between that wording and what the note does is collected in Remark \ref{c:note.angle.departures-registered-wording}.

\phantomsection\label{c:note.angle.setting-claim-kill-conditions-fixed-before-2}
\textbf{Claim G (the singularity paper's open route (a), as registered).} The Hawking angle follows from a spectrum condition with respect to a generator the ascent installs, within the dynamics the framework supplies: positivity of the energy of the boundary modes for such a generator selects a state whose one-sided marginals are boost-KMS at a finite $\beta$, and the Bisognano--Wichmann mechanism then fixes $\beta = 2\pi$ in rapidity.

\begin{itemize}
\item[K1] (the supplied class selects Boulware) If every joint-number-diagonal generator with a unique ground state has a number basis state as that ground state, then within the supplied class the spectrum condition selects a product state across every cut, whose one-sided marginals are pure: boost-KMS at $\beta = \infty$ only, the Boulware member of the dictionary. Claim G is dead within the supplied class. Test: $\nu = 2$ and $\nu = 3$, generic rational diagonal generators, and the degenerate case exhibited (a degenerate ground space can contain an entangled state, but then positivity selects nothing).
\item[K2] (a hop is what the route needs) If a generator with a hop across the cut, outside the supplied class, has a unique entangled ground state whose one-sided marginal is Gibbs for the number operator at a finite $\beta$, then the route survives exactly when the singularity paper's residual is answered with a hopping term: its open items (a) and (b) are one item. Test: two modes, one-particle sector, $H = \begin{psmallmatrix} 7 & -12 \\ -12 & 0\end{psmallmatrix}$ in the basis $\ket{10}, \ket{01}$; predicted ground state $\tfrac35\ket{10} + \tfrac45\ket{01}$, left marginal $\operatorname{diag}(\tfrac{16}{25}, \tfrac{9}{25})$, the sample point $e^{-\omega/T} = \tfrac{9}{16}$.
\item[K3] (the pairing is the hop with the sea filled) If the particle-hole map on the right mode carries the hop-type ground state $c\ket{10} + s\ket{01}$ to the pairing-type thermofield state $c\ket{00} + s\ket{11}$ and the hop generator to the pairing generator, then the hop with a filled sea on the right and the pairing from the vacuum are the same correlation in two vacua, and the residual ``differential or pairing term'' is one term seen from two vacua. Test: exact, on the two-mode module.
\item[K4] (the geometric half of Bisognano--Wichmann is in the algebra) If the boost rotor of the horizon plane at imaginary rapidity $i\pi$ acts on vectors as the reflection through the bifurcation surface, and if the translations of the conformal completion are the null rotations fixing $\nInf$, $T_{a} = 1 + \tfrac12\nInf a$ with $T_{a}X(x)T_{a}^{-1} = X(x + a)$, then the two geometric ingredients of the theorem, the wedge reflection and the translation group, are supplied by the ascent. Expected to hold; recorded as the positive residue, not as evidence for Claim G.
\item[K5] (the value) The factor $2\pi$ is the conclusion of Bisognano--Wichmann for a Wightman theory, or of Hislop--Longo for the half-line algebra of a chiral field, both statements about infinitely many modes with an analytic (spectrum) condition. If no finite mode model reproduces $2\pi$ from its own data, exact arithmetic reaches the necessity of a hop and not the value; the value stays imported. Expected to trigger.
\item[K6] (ledger) If the ascent's contribution is the vocabulary (translations, the boost plane, the rungs as modes) and no step of the argument needs the $(1,1)$ increment, the result is transport-absent. Expected.
\end{itemize}
Predictions, as registered: K1 triggers, the unique ground state of a diagonal generator being a number state with pure marginals, Gibbs at no finite rational $\beta$, and a degenerate diagonal generator admitting $c\ket{00} + s\ket{11}$ as a ground state but selecting no state (P1); K2 holds with the numbers stated, the ground state being odd, entangled across the cut, with right marginal $\operatorname{diag}(\tfrac{9}{25}, \tfrac{16}{25})$ (P2); K3 holds, the reflection carrying $\tfrac35\ket{10} + \tfrac45\ket{01}$ to $\tfrac35\ket{00} + \tfrac45\ket{11}$ up to sign and the hop Hamiltonian to a pairing Hamiltonian with the same spectrum (P3); K4 holds exactly, in $\Cl(1,3)$ for the reflection and in $\Cl(2,4)$ for the translations (P4); K5 triggers, no finite model supplying $2\pi$, and the note records the imports that do and where they apply, Wightman fields, the chiral half-line, bifurcate Killing horizons (P5); transport-absent (P6).

\section{The geometric half}

\begin{proposition}[The wedge reflection and the translations are in the ascent]\label{c:angle.wedge-reflection-translations-are-ascent}
In $\Cl(1,3)$ with $e_{t}^{2} = +1$ and $e_{r}^{2} = e_{\theta}^{2} = e_{\varphi}^{2} = -1$, let $B = e_{t}e_{r}$ be the horizon's boost plane, the boost member of the singularity paper's curvature-selected pair.
\begin{enumerate}
\item[(i)] $B^{2} = +1$; the rotor $R = \cosh\eta + B\sinh\eta$ at $(\cosh\eta, \sinh\eta) = (\tfrac54, \tfrac34)$ satisfies $R\widetilde{R} = 1$ and acts on $(e_{t}, e_{r})$ by the boost of rapidity $2\eta$, the matrix $\begin{psmallmatrix} 17/8 & -15/8 \\ -15/8 & 17/8\end{psmallmatrix}$, fixing $e_{\theta}$ and $e_{\varphi}$; it rescales the null generators $e_{t} \pm e_{r}$ of the plane by $e^{\mp 2\eta} = \tfrac14, 4$.
\item[(ii)] At boost rapidity $2\eta = i\pi$, that is $\eta = i\pi/2$, the rotor is $iB$, with $(iB)\widetilde{(iB)} = B^{2} = 1$, and its action $x \mapsto BxB$ is $-1$ on the boost plane and $+1$ transverse to it: the reflection through the bifurcation surface. It commutes with the boost and squares to the identity.
\item[(iii)] In the completion $\Cl(2,4)$, for $a$ in the base, $T_{a} = 1 + \tfrac12\nInf a$ is an even element with inverse $1 - \tfrac12\nInf a$, fixes $\nInf$, is the null rotation $R_{u}$ of Fact \ref{c:angle.imported-conformal-model-null-rotations} with $v = \nInf$ and $u = \tfrac12 a$, and carries $X(x)$ to $X(x + a)$. The base's boost rotor, read in the completion, fixes $\nO$ and $\nInf$ and carries $X(x)$ to $X(Rx\widetilde{R})$.
\end{enumerate}
\end{proposition}
\begin{proof}
All clauses are identities in the blade basis, checked by the script (Parts A and B) on rational vectors. For (ii): $\cosh(i\pi/2) + B\sinh(i\pi/2) = iB$, whose reverse is $-iB$, so the twisted action is $x \mapsto (iB)x(-iB) = BxB$; $B$ anticommutes with $e_{t}$ and $e_{r}$ and commutes with $e_{\theta}$, $e_{\varphi}$, and $B^{2} = 1$. For (iii): $(\nInf a)^{2} = -\nInf^{2}a^{2} = 0$ since $a$ anticommutes with $\nInf$, so $T_{a}$ is invertible with the stated inverse and is the null rotation $R_{u}$ of \ref{c:angle.imported-conformal-model-null-rotations} with $u = \tfrac12 a$. Expanding $T_{a}yT_{a}^{-1} = y + \tfrac12[\nInf a, y] - \tfrac14\nInf a\,y\,\nInf a$ and using $\nInf^{2} = 0$, $\nInf a = -a\nInf$ and $\nO\nInf + \nInf\nO = -2$: for $y = \nInf$ every correction vanishes; for $y = x$ in the base the commutator term gives $(x \cdot a)\nInf$ and the last term vanishes; for $y = \nO$ the commutator term gives $a$ and the last term $\tfrac12 Q(a)\nInf$. Summing with the coefficients of $X(x)$ gives $\nO + (x + a) + \tfrac12(Q(x) + 2x\cdot a + Q(a))\nInf = X(x + a)$, in the sign convention $\nO \cdot \nInf = -1$.

\end{proof}

\begin{remark}[What this supplies, and what it does not]\label{c:note.angle.what-supplies-what-it-does-not}
The Bisognano--Wichmann theorem has two halves. The geometric half says that the boost at rapidity $i\pi$ is the reflection exchanging a wedge with its causal complement, and that the translations along a null generator of the horizon, $a \in \R_{\geq 0}(e_{t} + e_{r})$, move the wedge into itself and are rescaled by the boost, $R(e_{t} + e_{r})\widetilde{R} = e^{-2\eta}(e_{t} + e_{r})$; the base supplies the boost plane and the ascent supplies the translations as the null rotations of the point at infinity, exactly as the singularity paper's remark anticipated. The selective half says that a state with positive energy for the translations has, for the wedge algebra, a modular group equal to the boost at $2\pi$; that is a statement about the state, and it is the subject of the next section.
\end{remark}

\section{The selective half, in the algebra}

\begin{proposition}[Within the supplied class the spectrum condition selects Boulware, or nothing]\label{c:angle.within-supplied-class-spectrum-condition-selects}
Let $H$ be a generator in the supplied class on $\nu$ modes, a polynomial in the number operators, and let a cut divide the modes into two sides.
\begin{enumerate}
\item[(i)] $H$ is diagonal in the joint number basis. If its ground state is unique, that state is a number basis state: a product across the cut, with a pure marginal on each side, an eigenvector of every $N_{k}$ and hence of every boost generator $\sum_{k}\omega_{k}(N_{k} - \tfrac12)$, and Gibbs on a side at no finite $\beta$: a Boulware-type state, a product with pure marginals.
\item[(ii)] If the ground state is not unique, the ground space may contain entangled states, but the spectrum condition selects none of them: for $H = (N_{0} - N_{1})^{2} = N_{0} + N_{1} - 2N_{0}N_{1}$ on the pair, the ground space is $\operatorname{span}\{\ket{00}, \ket{11}\}$, which contains the thermofield pair $\tfrac45\ket{00} + \tfrac35\ket{11}$ and the product state $\ket{00}$ alike.
\end{enumerate}
Within the supplied class, therefore, positivity of the energy selects a Boulware-type state or selects no state, and Claim G is dead at the state-selection step.
\end{proposition}
\begin{proof}
(i) By Theorem \ref{c:singularity.class-relative-no-go-no-supplied} the supplied class is ultralocal, every generator in it a multiplication operator in the joint number basis, so a polynomial in the commuting projectors $N_{k}$. A polynomial in commuting projectors is diagonal in their joint eigenbasis; a diagonal matrix with a unique minimum has that basis vector as its unique ground state. A number basis state is a product of single-mode number states, its marginal on any set of modes is the corresponding number state, pure, and it is an eigenvector of each $N_{k}$. A Gibbs state $e^{-\beta K}/Z$ at finite $\beta$ for a single-mode boost $K = \omega(N - \tfrac12)$ is $\operatorname{diag}(1, e^{-\beta\omega})/Z$, of rank two, and a pure marginal has rank one. The script draws random rational diagonal generators on two and three modes with every cut, and for each draw with a unique ground state (at least three such draws in all) confirms the product form, the rank-one idempotent marginals on both sides, the eigenvector property, and, for a single-mode side, the absence of a Gibbs parameter. (ii) The stated $H$ has eigenvalues $0, 0, 1, 1$ on $\ket{00}, \ket{11}, \ket{10}, \ket{01}$; the script checks the ground space and that both named states are annihilated by $H$.

\end{proof}

\begin{proposition}[A hop across the cut]\label{c:angle.hop-across-cut}
On two modes, let $H = 7N_{0} - 12(a^{\dagger}_{0}a_{1} + a^{\dagger}_{1}a_{0})$, the second quantisation of the one-particle operator $\begin{psmallmatrix} 7 & -12 \\ -12 & 0\end{psmallmatrix}$, which is not in the supplied class.
\begin{enumerate}
\item[(i)] $H$ is Hermitian and conserves the total number; its spectrum is $0$ on $\ket{00}$, $7$ on $\ket{11}$, and $16, -9$ on the one-particle sector, so its ground state is unique: $\psi = \tfrac35\ket{10} + \tfrac45\ket{01}$ with eigenvalue $-9$.
\item[(ii)] $\psi$ is an odd vector, hence of definite parity, so the state it defines is even (it vanishes on the odd sector) and is physical under the univalence rule, and entangled across the cut: its left marginal is $\operatorname{diag}(\tfrac{16}{25}, \tfrac{9}{25})$, the Gibbs state for $N_{0}$ at $e^{-\omega/T} = \tfrac{9}{16}$, the singularity paper's sample point, and its right marginal is $\operatorname{diag}(\tfrac{9}{25}, \tfrac{16}{25})$.
\item[(iii)] $\psi$ is an eigenvector of $N_{0} + N_{1}$ and not of $N_{0} - N_{1}$: it is invariant under the boost flow in which the two modes have equal frequency, and not under the number-difference boost of the vacuum form.
\end{enumerate}
\end{proposition}
\begin{proof}
The one-particle block has trace $7$ and determinant $-144$, eigenvalues $16$ and $-9$, and the eigenvector for $-9$ satisfies $16u = 12v$, giving $(u, v) = (\tfrac35, \tfrac45)$. The script checks the block, the sector energies, the eigenvector, the norm, the parity, both marginals by restriction to the mode's algebra, and the two eigenvector statements. The generator lies outside the supplied class because every supplied generator commutes with every $N_{x}$ (Theorem \ref{c:singularity.class-relative-no-go-no-supplied}), while $[N_{0}, a^{\dagger}_{0}a_{1} + a^{\dagger}_{1}a_{0}] = a^{\dagger}_{0}a_{1} - a^{\dagger}_{1}a_{0} \neq 0$.

\end{proof}

\begin{proposition}[The hop with the sea filled is the pairing]\label{c:angle.hop-sea-filled-pairing}
Let $V = (a_{1} + a^{\dagger}_{1})\Gamma_{0}$ on the pair, with $\Gamma_{0} = (-1)^{N_{0}}$ and $\ket{11} = a^{\dagger}_{0}a^{\dagger}_{1}\ket{00}$. Then $V$ is unitary and Hermitian, $V^{2} = 1$, conjugation by $V$ fixes $a_{0}$ and exchanges $a_{1}$ with $a^{\dagger}_{1}$ (the particle-hole automorphism of the right mode), and:
\begin{enumerate}
\item[(i)] $V\psi = \tfrac45\ket{00} + \tfrac35\ket{11}$, exactly the singularity paper's thermofield pair at the sample point, an even state with the same left marginal $\operatorname{diag}(\tfrac{16}{25}, \tfrac{9}{25})$;
\item[(ii)] $VHV^{-1} = 7N_{0} - 12(a^{\dagger}_{0}a^{\dagger}_{1} + a_{1}a_{0})$, a pairing generator with the same spectrum, whose unique ground state is $V\psi$;
\item[(iii)] $V(N_{0} + N_{1} - 1)V^{-1} = N_{0} - N_{1}$: the equal-frequency boost of the hop form becomes the number-difference boost of the vacuum form.
\end{enumerate}
\end{proposition}
\begin{proof}
$(a_{1} + a^{\dagger}_{1})$ is a Majorana operator, Hermitian with square one, and conjugation by it sends $a_{1}$ to $a^{\dagger}_{1}$ and $a_{0}$ to $-a_{0}$; the factor $\Gamma_{0}$ restores the sign of $a_{0}$ and commutes with the right mode. The script checks the unitarity, the automorphism on the four ladder operators, the image of $\psi$, its parity and marginal, the conjugated generator against the pairing generator entry by entry, its ground state, and the conjugated boost. Here $\psi$ and $H$ are the state and generator of Proposition \ref{c:angle.hop-across-cut}; conjugation by the unitary $V$ preserves the spectrum, so $VHV^{-1}$ has that proposition's spectrum and its unique ground state is the image $V\psi$ of the unique ground state there.

\end{proof}

\begin{remark}[One term in two vacua]\label{c:note.angle.one-term-two-vacua}
The singularity paper's residual question asks whether the physical dynamics on the boundary carries a differential term or a pairing term, the two kinds of generator its supplied class lacks. Proposition \ref{c:angle.hop-sea-filled-pairing} says the question has one answer, not two: a hop between a mode and its partner, with the partner's sea filled, is a pairing between the mode and the partner's hole, and the ground state of the one is the thermofield pair of the other. That a number-conserving generator applied to the vacuum keeps the pair in its vacuum, while a hop applied to a one-particle state produces the thermofield marginal, is the same fact read in the two vacua: what the hop needs is the filled sea the s-wave array of Step 4e supplies, and what the pairing needs is the vacuum. The particle-hole map $V$ is the Fock form of a symmetry of the ascent itself: in the two-rung ascent $\Cl(2,2)$ with ladder pairs $f_{k} = \tfrac12(e_{k} + e_{k+2})$, $f_{k}^{*} = \tfrac12(e_{k} - e_{k+2})$, the reflection $e_{4} \mapsto -e_{4}$ is an automorphism that fixes $f_{1}$, carries $f_{2}$ to $f_{2}^{*}$, and carries the hop $f_{1}^{*}f_{2}$ to the pairing $f_{1}^{*}f_{2}^{*}$ (script, Part E).
\end{remark}

\begin{proposition}[The value is a parameter of the dynamics]\label{c:angle.value-parameter-dynamics}
For every rational point $(c, s)$ of the unit circle with $cs \neq 0$ and every rational $t$ with $tcs > 0$, the hop generator $H = \varepsilon N_{0} - t(a^{\dagger}_{0}a_{1} + a^{\dagger}_{1}a_{0})$ with $\varepsilon/t = (s^{2} - c^{2})/(cs)$ has $c\ket{10} + s\ket{01}$ as its unique ground state, with eigenvalue $-tc/s$, and left marginal Gibbs for $N_{0}$ at $q = c^{2}/s^{2}$. The Gibbs parameter of the marginal, and with it $\beta\omega = -\log q$, is determined by the ratio $\varepsilon/t$ of the couplings together with the sign of $t$ (by the ratio alone for $t > 0$): $c/s$ is the root of $r^{2} + (\varepsilon/t)r - 1 = 0$ with the sign of $t$.
\end{proposition}
\begin{proof}
The one-particle block $\begin{psmallmatrix} \varepsilon & -t \\ -t & 0\end{psmallmatrix}$ has $(c, s)$ as an eigenvector with eigenvalue $\lambda = -tc/s$ exactly when $\varepsilon c - ts = \lambda c$, i.e.\ $\varepsilon = t(s^{2} - c^{2})/(cs)$; the other one-particle eigenvalue is $\varepsilon - \lambda$, and the three gaps, to it and to the sector energies $0$ and $\varepsilon$, are $ts/c + tc/s$, $tc/s$ and $ts/c$, all positive exactly when $tcs > 0$. The script confirms the eigenvector, the gaps and the marginal at seven points, $(\tfrac35, \tfrac45)$, $(\tfrac{5}{13}, \tfrac{12}{13})$, $(\tfrac{8}{17}, \tfrac{15}{17})$, $(\tfrac{7}{25}, \tfrac{24}{25})$, $(\tfrac{24}{25}, \tfrac{7}{25})$ and $(-\tfrac35, -\tfrac45)$ at $t = 60$ and $(\tfrac35, -\tfrac45)$ at $t = -60$. Given $\varepsilon/t$, the ratio $c/s$ is the root of $r^{2} + (\varepsilon/t)r - 1 = 0$ with the sign of $t$ (the two roots have product $-1$) and $q$ is its square.

\end{proof}

\begin{remark}[Where $2\pi$ lives]\label{c:note.angle.where-2-lives}
Proposition \ref{c:angle.value-parameter-dynamics} is the finite-model form of the singularity paper's rate rigidity: a hop selects an entangled state and a finite $\beta$, but the value of $\beta$ is whatever the couplings make it. The number $2\pi$ enters only through Fact \ref{c:angle.imported-theorems-deliver-their-hypotheses}, whose abstract form needs a half-sided translation group with positive generator on the mode algebra: infinitely many modes, the translation acting by shifting them, the state its ground state, and the algebra of one side mapped into itself by the translations. A ${\mathbb{Z}}$-array with a hopping term and a filled sea is the discrete shadow of that structure; the theorem itself is a statement about its continuum limit and its analytic continuation, which exact arithmetic does not reach. The note reaches the necessity of the hop and stops there.
\end{remark}

\section{What this does to the singularity paper, and what it does not do}

\phantomsection\label{c:note.angle.setting-what-does-singularity-paper-what}
\textbf{Verdict.} K1 triggered: within the supplied class the spectrum condition selects a Boulware-type state, a product with pure marginals, or nothing (Proposition \ref{c:angle.within-supplied-class-spectrum-condition-selects}), so Claim G is dead at the state-selection step, as registered. K2, K3 and K4 held (Propositions \ref{c:angle.hop-across-cut}, \ref{c:angle.hop-sea-filled-pairing}, \ref{c:angle.wedge-reflection-translations-are-ascent}). K5 triggered (Proposition \ref{c:angle.value-parameter-dynamics}, Remark \ref{c:note.angle.where-2-lives}). K6: the base supplies the boost plane, the ascent supplies the translations and the modes, and no step of the argument uses the $(1,1)$ increment; transport-absent. Of the predictions, P1, P2, P4, P5 and P6 held as registered; P3 held with its amplitudes corrected as Remark \ref{c:note.angle.departures-registered-wording} records.

\phantomsection\label{c:note.angle.setting-what-does-singularity-paper-what-2}
\textbf{One item, not two.} The singularity paper leaves two things open at the horizon: (a) whether (R) can be traded for a spectrum condition in the ascent's own algebra, and (b) whether the physical dynamics on the boundary carries a differential or a pairing term. Propositions \ref{c:angle.within-supplied-class-spectrum-condition-selects} to \ref{c:angle.hop-sea-filled-pairing} join them. Within the supplied class the spectrum condition selects a product state or nothing; a hop across the cut can select an entangled state, and the example does, at the sample point; a hop with the sea filled and a pairing from the vacuum are one term in two vacua. So (a) is answered affirmatively only if the boundary dynamics carries a generator outside the supplied class, which is the residual (b) declares; the paper's residual names the two kinds, and the example uses a hop across the cut. The trade on offer is conditional: (R), regularity at the fixed locus, can be replaced by the premise that the boundary dynamics carries a kinetic term with positive energy under the translations the ascent installs, together with the remaining hypotheses of Fact \ref{c:angle.imported-theorems-deliver-their-hypotheses}, the half-sided inclusion of the one-sided algebra, the cyclic and separating property of the state, an action of the installed translations on the modes themselves, and the identification of the modular flow with the Killing boost, none of which this note verifies for the boundary modes; under those, the value $2\pi$ is the imported theorem's. That is one physical premise, a kinetic term, in place of one regularity premise, plus a theorem with hypotheses still to be met. It is not a derivation of the temperature; it is a relocation of the import to a statement about the dynamics rather than about the state, which is the relocation the paper asked for.

\begin{remark}[Departures from the registered wording]\label{c:note.angle.departures-registered-wording}
Three. (1) K3 and P3 as registered wrote the image of $c\ket{10} + s\ket{01}$ under the particle-hole map as $c\ket{00} + s\ket{11}$, up to sign. The image is $s\ket{00} + c\ket{11}$: $V$ inverts the right occupation, sending $\ket{01}$ to $\ket{00}$ and $\ket{10}$ to $\ket{11}$, and fixes $a_{0}$, so the left marginal is preserved and the amplitudes exchange places; at the sample point the image is the singularity paper's pair $\tfrac45\ket{00} + \tfrac35\ket{11}$ rather than its transpose, and the sign is that of the convention $\ket{11} = a^{\dagger}_{0}a^{\dagger}_{1}\ket{00}$. K3 is read as held because its substance, that the image is the thermofield pair with the same marginal and the hop generator becomes a pairing generator, holds. (2) K5 as registered attributed the chiral half-line case to Hislop--Longo; the note cites Hislop--Longo for the free massless field and the conformal-net theorem for the chiral case, and states the abstract form as Borchers' theorem, which is the form the route needs. (3) Proposition \ref{c:angle.wedge-reflection-translations-are-ascent}(i)'s rescaling of the null generators and Remark \ref{c:note.angle.one-term-two-vacua}'s exchange in the ascent are outside the registration.
\end{remark}

\phantomsection\label{c:note.angle.setting-what-does-singularity-paper-what-3}
\textbf{Prior art, and the two ledgers.} The Bisognano--Wichmann theorem, its abstract form and its bifurcate-horizon form are imports; the Boulware, Hartle--Hawking and Unruh states are imports; the particle-hole symmetry of a fermionic mode is elementary. What the note proves for itself is the geometric half in the ascent's terms, the Boulware selection within the supplied class, the hop example at the sample point, the identification of the hop with the pairing under the particle-hole map, and the parametrisation of the Gibbs parameter by the couplings. Under the two-ledger discipline the result is transport-absent: the base supplies the boost plane, the $(1,1)$ increment supplies the modes and the translations, all of which any fermionic mode algebra over a spacetime with a bifurcate Killing horizon supplies, and the argument uses none of them as a step \cite{ButlerT17}.

\phantomsection\label{c:note.angle.setting-what-does-singularity-paper-what-4}
\textbf{What it does not do.} It does not derive $2\pi$; Remark \ref{c:note.angle.where-2-lives} says where the number lives and why the finite model cannot reach it. It does not say whether nature's dynamics on the singularity paper's boundary carries a hop; that remains the residual, now stated once. It does not move the modes: the singularity paper attaches them to the exceptional set and calibrates at the horizon, and Proposition 7.25 of that paper stands. And it does not touch the vacuum of the spinor fibre at a point, the null direction a Witt decomposition of the fibre selects: the selection rule that exists, positivity of energy, is a property of a generator acting across the cut, and no datum at one point constrains a correlation between two.

\section{Position in the series}

\phantomsection\label{c:note.angle.setting-position-series}
This document carries no T-number: it is a companion note, deposited with T5 \cite{ButlerT5} and read with T1 \cite{ButlerT1}, and its notation is that of the primer T0 \cite{ButlerT0}. From T5 it takes the state selection of Step 4d (Lemma 7.2, Theorems 7.3 and 7.16, Propositions 7.18 and 7.25, Premise 7.19, Remarks 7.26 and 7.33), the supplied class and its class-relative no-go (Premise 7.30, Theorem C.8, Remark C.9), the s-wave array with its sea (Premise 8.1), and the surface gravity (Proposition A.2). From T1 it takes the conformal model and the null rotations (Section 4.3, Lemma 4.6, Theorem 4.10). It sends back to T5 the merger of its open items (a) and (b) and the trade of the section ``What this does to the singularity paper, and what it does not do''. The capstone T17 \cite{ButlerT17} carries the two-ledger grading used in the section ``What this does to the singularity paper, and what it does not do''.

\section{Machine verification}

\phantomsection\label{c:note.angle.setting-machine-verification}
Every computational identity asserted in this note has been verified in exact arithmetic by the script \texttt{verify\_T05\_angle.py} \cite{ButlerVerifyT05angle}, which imports the series' shared library \texttt{GA.py}: its blade engine for the geometric parts and its Gaussian rationals for the Fock parts. The reduced state of a subsystem is computed as the restriction of the joint state to the subsystem's algebra, the fermionic reduced state, so that no tensor-product identification is needed. No floating point is used anywhere; the script exits nonzero on any failure. Eighteen checks, all passing, in six parts. The pre-registration is deposited beside the script \cite{ButlerPreregT05angle}.

\emph{Part A: the boost plane.} $B = e_{t}e_{r}$ is a bivector of square $+1$; the rotor at $(\tfrac54, \tfrac34)$ has $R\widetilde{R} = 1$ and the stated boost matrix on $(e_{t}, e_{r})$, fixing the transverse vectors; $x \mapsto BxB$ is $-1$ on the plane and $+1$ transverse, with $B\widetilde{B} = -1$ so that $(iB)\widetilde{(iB)} = 1$; the reflection commutes with the boost and is an involution on vectors; the boost rescales $e_{t} + e_{r}$ by $\tfrac14$ and $e_{t} - e_{r}$ by $4$.

\emph{Part B: the translations.} $\nO$, $\nInf$ null with $\nO\cdot\nInf = -1$; $X(x)$ null and $X(x)\cdot X(y) = -\tfrac12 Q(x - y)$ on rational points; $T_{a}$ even, with the stated inverse, fixing $\nInf$, of the null-rotation form with $u = \tfrac12 a$ orthogonal to $\nInf$; $T_{a}X(x)T_{a}^{-1} = X(x + a)$ and $T_{a}X(x - a)T_{a}^{-1} = X(x)$; the base's boost rotor fixes $\nO$ and $\nInf$ and commutes with the lift.

\emph{Part C: the supplied class.} Random rational diagonal generators on two and three modes with every cut: diagonal in the number basis; for each unique ground state (at least three such draws asserted), a number state with rank-one idempotent marginals on both sides, an eigenvector of every $N_{k}$, and, for a single-mode side, no Gibbs parameter; the degenerate generator $(N_{0} - N_{1})^{2}$ with ground space $\operatorname{span}\{\ket{00}, \ket{11}\}$ annihilating both the thermofield pair and $\ket{00}$.

\emph{Part D: the hop.} $H$ Hermitian; the one-particle block $\begin{psmallmatrix} 7 & -12 \\ -12 & 0\end{psmallmatrix}$; the eigenvector $\tfrac35\ket{10} + \tfrac45\ket{01}$ at $-9$; sector energies $0$ and $7$ with no off-diagonal entries; trace and determinant giving $16$ and $-9$; the ground state odd, of norm one, with both marginals of rank two, the left $\operatorname{diag}(\tfrac{16}{25}, \tfrac{9}{25})$ at $q = \tfrac{9}{16}$ and the right $\operatorname{diag}(\tfrac{9}{25}, \tfrac{16}{25})$; an eigenvector of $N_{0} + N_{1}$ and not of $N_{0} - N_{1}$.

\emph{Part E: the particle-hole map.} $V$ unitary, Hermitian, an involution, fixing $a_{0}$ and exchanging $a_{1}$ with $a^{\dagger}_{1}$; $V\psi = \tfrac45\ket{00} + \tfrac35\ket{11}$, even, with left marginal $\operatorname{diag}(\tfrac{16}{25}, \tfrac{9}{25})$; $VHV^{-1}$ equal to the pairing generator entry by entry, with $V\psi$ its eigenvector at $-9$; $V(N_{0} + N_{1} - 1)V^{-1} = N_{0} - N_{1}$; in $\Cl(2,2)$ the reflection $e_{4} \mapsto -e_{4}$ is an automorphism (random elements) fixing $f_{1}$ and carrying $f_{2}$ to $f_{2}^{*}$ and the hop to the pairing.

\emph{Part F: the parameter.} At the seven rational points, with $tcs > 0$, the hop generator with $\varepsilon/t = (s^{2} - c^{2})/(cs)$ has $c\ket{10} + s\ket{01}$ as an eigenvector at $-tc/s$, the three gaps equal to $ts/c + tc/s$, $tc/s$ and $ts/c$ and positive, and left marginal Gibbs at $q = c^{2}/s^{2}$.

What the script does not check: Facts \ref{c:angle.imported-singularity-papers-state-selection} to \ref{c:angle.imported-theorems-deliver-their-hypotheses} (imports); the general-$\nu$ clause of Proposition \ref{c:angle.within-supplied-class-spectrum-condition-selects}(i) beyond $\nu = 3$ (the proof is a statement about diagonal matrices); the uniqueness clause of Proposition \ref{c:angle.value-parameter-dynamics} beyond the seven points (the proof is the two-by-two computation); the Gibbs clause of Proposition \ref{c:angle.within-supplied-class-spectrum-condition-selects}(i) for sides of more than one mode (the rank argument); and the continuation to imaginary rapidity in Proposition \ref{c:angle.wedge-reflection-translations-are-ascent}(ii), of which the script checks the endpoint $iB$ and its action.

\section{Results used from elsewhere}
The proofs above use the following results of the series, stated briefly; each is proved in the paper named.

\begin{fact}[The class-relative no-go: no supplied generator hops]\label{c:singularity.class-relative-no-go-no-supplied}
Given function-linearity of the tame constructions, no generator in the supplied class moves occupation between points of $E$, while pairing generators lie outside the class and are not excluded.
\emph{Source:} T5, Thm C.8 (\href{https://doi.org/10.5281/zenodo.22721790}{doi:10.5281/zenodo.22721790}).
\end{fact}

\begin{fact}[$\kappa = 1/4M$, computed]\label{c:singularity.computed}
For the Schwarzschild field on flat spacetime the redshifted acceleration of static observers, $v = K/\sqrt f$, is $f'(r)/2$, so the surface gravity is $\kappa = 1/4M$.
\emph{Source:} T5, Prop A.2 (\href{https://doi.org/10.5281/zenodo.22721790}{doi:10.5281/zenodo.22721790}).
\end{fact}

\begin{fact}[The Euclidean section, and its period]\label{c:singularity.euclidean-section-period}
The Euclidean section of the Schwarzschild field is smooth at its tip, the Euclidean counterpart of the bifurcation sphere, if and only if imaginary Killing time has period $p = 2\pi/\kappa = 8\pi M$; any other period leaves a conical point of total angle $\kappa p$.
\emph{Source:} T5, Prop 7.18 (\href{https://doi.org/10.5281/zenodo.22721790}{doi:10.5281/zenodo.22721790}).
\end{fact}

\begin{fact}[Modular descent]\label{c:singularity.modular-descent}
A pure, boost-invariant, entangling state of the horizon pair descends to a KMS state of the phase flow $\operatorname{Ad}(e^{it\sigma/2})$ at a free $\beta$, and under the geometric normalisation (R) to the Fermi-Dirac state at exactly $\TH = \kappa/2\pi = 1/8\pi M$.
\emph{Source:} T5, Thm 7.3 (\href{https://doi.org/10.5281/zenodo.22721790}{doi:10.5281/zenodo.22721790}).
\end{fact}

\begin{fact}[Rate rigidity: the temperature is not algebra intrinsic]\label{c:singularity.rate-rigidity-temperature-not-algebra-intrinsic}
Automorphisms of the single-mode algebra that intertwine the boost's phase flow $\operatorname{Ad}(e^{itK})$ move its KMS states only by $\beta \mapsto \pm\beta$, so the value of the temperature depends on a parametrisation of the flow that the algebra does not supply.
\emph{Source:} T5, Thm 7.16 (\href{https://doi.org/10.5281/zenodo.22721790}{doi:10.5281/zenodo.22721790}).
\end{fact}

\begin{fact}[Schmidt forcing, proved outright]\label{c:singularity.schmidt-forcing-proved-outright}
A unit vector on the doubled mode that is boost invariant up to phase and cyclic and separating is $\cos\theta\,|00\rangle + \sin\theta\,|11\rangle$ with both coefficients nonzero, so its Schmidt basis is the number basis.
\emph{Source:} T5, Lem 7.2 (\href{https://doi.org/10.5281/zenodo.22721790}{doi:10.5281/zenodo.22721790}).
\end{fact}

\begin{fact}[The group of null rotations]\label{c:transport.group-null-rotations}
The null rotations $R_u = 1 + vu$ with $u \in U$ form an abelian subgroup of $\Spin^0(p,q)$ that fixes $v$ and is isomorphic to $(U, +)$, equivariantly for the rotors of $U$.
\emph{Source:} T1, Lem 4.6 (\href{https://doi.org/10.5281/zenodo.21743863}{doi:10.5281/zenodo.21743863}).
\end{fact}

\begin{fact}[Totality]\label{c:transport.totality}
On the quadric of the completion, $\iota[N] = [e_+ N e_+]$ is a total continuous involution that equals the reciprocal $x \mapsto x/Q(x)$ on invertible vectors and sends each $x_0$ with $Q(x_0) = 0$ to $[x_0 + \frac12 \nInf]$, and it is the only such continuous extension.
\emph{Source:} T1, Thm 4.10 (\href{https://doi.org/10.5281/zenodo.21743863}{doi:10.5281/zenodo.21743863}).
\end{fact}

\begin{thebibliography}{99}
\bibitem{ButlerT5} L.~Butler,
\emph{Division by Zero at the Gravitational Singularity: The Transport Principle Applied to the Black Hole, and the Information Paradox as a Category Error},
Transport Papers T5 (2026), doi:\href{https://doi.org/10.5281/zenodo.22721790}{10.5281/zenodo.22721790}.
\bibitem{ButlerT1} L.~Butler,
\emph{The Transport Principle: Total Division in Clifford Algebras via Minimal Completion},
Transport Papers T1 (2026).
\bibitem{BW1975} J.~J.~Bisognano and E.~H.~Wichmann,
On the duality condition for a Hermitian scalar field,
\emph{J.\ Math.\ Phys.} \textbf{16} (1975), 985--1007.
\bibitem{BW1976} J.~J.~Bisognano and E.~H.~Wichmann,
On the duality condition for quantum fields,
\emph{J.\ Math.\ Phys.} \textbf{17} (1976), 303--321.
\bibitem{Borchers1992} H.-J.~Borchers,
The CPT-theorem in two-dimensional theories of local observables,
\emph{Commun.\ Math.\ Phys.} \textbf{143} (1992), 315--332.
\bibitem{HislopLongo1982} P.~D.~Hislop and R.~Longo,
Modular structure of the local algebras associated with the free massless scalar field theory,
\emph{Commun.\ Math.\ Phys.} \textbf{84} (1982), 71--85.
\bibitem{BGL1993} R.~Brunetti, D.~Guido and R.~Longo,
Modular structure and duality in conformal quantum field theory,
\emph{Commun.\ Math.\ Phys.} \textbf{156} (1993), 201--219.
\bibitem{Sewell1982} G.~L.~Sewell,
Quantum fields on manifolds: PCT and gravitationally induced thermal states,
\emph{Ann.\ Phys.} \textbf{141} (1982), 201--224.
\bibitem{KayWald1991} B.~S.~Kay and R.~M.~Wald,
Theorems on the uniqueness and thermal properties of stationary, nonsingular, quasifree states on spacetimes with a bifurcate Killing horizon,
\emph{Phys.\ Rep.} \textbf{207} (1991), 49--136.
\bibitem{Unruh1976} W.~G.~Unruh,
Notes on black-hole evaporation,
\emph{Phys.\ Rev.\ D} \textbf{14} (1976), 870--892.
\bibitem{ButlerPreregT05angle} L.~Butler,
\emph{preregistration\_T05\_angle.md: the kill conditions and predictions of ``What Fixes the Angle'', as registered before computation},
deposited with this note (2026).
\bibitem{ButlerT17} L.~Butler,
\emph{Capstone: The Transport Programme, Method, Status, and the Falsifiable Ledger},
Transport Papers T17 (2026).
\bibitem{ButlerT0} L.~Butler,
\emph{An Introduction to Geometric Algebra: A Primer for the Transport Papers},
Transport Papers T0 (2026).
\bibitem{ButlerVerifyT05angle} L.~Butler,
\emph{verify\_T05\_angle.py: exact-arithmetic verification of the computational identities of ``What Fixes the Angle''},
Python script, deposited with this note alongside the shared library GA.py (2026).
\end{thebibliography}
\end{document}